% Isolated research candidate, 2026-08-31. No independent audit, production
% acceptance, live input, or inventory closure is asserted.
Let \(k\) be algebraically closed and trivially valued, let
\(K=k((t))\), \(|t|=\mathrm e^{-1}\), and let \(X/k\) be smooth,
integral and projective with ample line bundle \(L\). Put
\(S=X\times\mathbf P^1_k\), with projections \(p,q\), base coordinates
\([T:U]\), and \(t=T/U\). If \(\eta=\operatorname{ord}_{t=0}\), its
completed residue field is \(K\), and its full analytic fibre is
\(F=q^{-1}(\eta)=X_K^{\mathrm{an}}\). This is not a Gauss-section image.

Use trivial metrics over \(k\), and the constant-model metric on \(L_K\)
from \(X\times\operatorname{Spec}k[[t]]\). A real-weighted FS potential
on \(L_K\) means
\begin{equation}\label{eq:p081-uniform-fs-definition}
 f=\frac1m\max_j\{\log|s_j|_0+\lambda_j\},\qquad
 s_j\in H^0(X_K,mL_K),\quad\lambda_j\in\mathbf R,
\end{equation}
with \(m\geq1\) and a finite section system having no common zero.
This is the metric convention of \cite[Definition~5.2]{BE2021}.
Rational FS means that the weights can be chosen rational. Over \(k\),
\(H_{\mathbf R}\) and \(H_{\mathbf Q}\) have exactly the meaning of
\cite[Definition~2.18]{BJ22}.

\begin{proposition}\label{prop:p081-uniform-fibre-fs-extension}
Fix \(M\geq0\) and an integer \(C\geq\max\{1,M\}\). Every FS
potential \(f\) in \([-M,0]\) is the exact restriction of
\[
 \Phi\in H_{\mathbf R}(A),\qquad
 A=p^*L\otimes q^*\mathcal O(C),
\]
such that \(\Phi|_F=f\) and \(-C\leq\Phi\leq1\) on
\(S^{\mathrm{an}}\). It has a section representation with all
normalized weights in \([-C,1)\). Rational inputs have rational
extensions. The fixed \(C\) is independent of denominator, section
count, coefficients and accuracy; the error is zero.

For a fixed FS reference \(\rho\in H_{\mathbf R}(L)\), choose
integers \(a,b\geq0\) with \(-b\leq\rho\leq a\). If an FS metric
on \(L_K\) has relative potential \(-M\leq\phi\leq0\) with respect
to \(\rho_K\), any integer \(C\geq\max\{1,M+a+b\}\) works: an FS
metric on \(A\) has relative potential \(H\), with respect to the
product reference \(p^*\rho\) and the trivial base metric, satisfying
\(H|_F=\phi\) and \(-C\leq H\leq a+b+1\).
\end{proposition}

Define \(\operatorname{PSH}_{\mathrm{FS}}(L_K)\) to consist of the
functions, relative to the constant-model reference, that are not
identically \(-\infty\) and are pointwise decreasing-net limits of FS
potentials on this fixed \(L_K\). Thus it is the nontrivial part of the
metric class in \cite[Definition~7.2]{BE2021}, expressed as relative
potentials. No identification with a different psh definition is implicit.

\begin{theorem}\label{thm:p081-constant-laurent-envelope-transfer-candidate}
Assume the BJ envelope property on \(S\) in every ample class
\(c_1(p^*L\otimes q^*\mathcal O(C))\), \(C\in\mathbf Z_{\geq1}\).
Then every nonempty family \((u_i)\subset
\operatorname{PSH}_{\mathrm{FS}}(L_K)\) bounded above by one real
constant satisfies
\[
                  (\sup_i u_i)^{*F}
                       \in\operatorname{PSH}_{\mathrm{FS}}(L_K).
\]
For FS inputs \(f_i\in[-M,0]\), the extensions above with one common
\(C\) satisfy the exact identity
\begin{equation}\label{eq:p081-uniform-usc-fibre}
              (\sup_i\Phi_i)^{*S}|_F=(\sup_i f_i)^{*F}.
\end{equation}
The stars use the indicated ambient or fibre topology.
\end{theorem}

The supplied unibranch-envelope paper asserts the hypothesis, since
\(S\) is smooth projective. The new bridge is candidate-only, pending
independent audit and production verification. Its scope is constant
polarized varieties over \(k((t))\), not arbitrary complete fields,
nonconstant varieties, arbitrary numerical classes on \(X_K\), or CLD
comparison. Continuous-obstacle continuity is neither assumed nor claimed.
